\documentclass[11pt,a4paper]{article}
\usepackage[margin=2.5cm]{geometry}
\usepackage{graphicx}
\usepackage{amsmath,amssymb}
\usepackage{booktabs}
\usepackage{hyperref}
\usepackage{xcolor}
\usepackage{float}
\usepackage{tabularx}
\usepackage{enumitem}

\tolerance=1000
\emergencystretch=3em
\hfuzz=2pt

\definecolor{success}{HTML}{2E7D32}
\definecolor{failure}{HTML}{C62828}
\definecolor{partial}{HTML}{F57F17}

\title{Boltzmann Thermodynamics of Attention:\\Level 1 Experiment Report}
\author{Intermediate Results --- \today}
\date{}

\begin{document}
\maketitle

\begin{abstract}
We test whether the formal equivalence between softmax attention and the Boltzmann distribution yields practical value for KV cache quantization.
Level~1 experiments measure (1)~effective temperature $T_\mathrm{eff}$ monotonicity across layers and (2)~whether the Boltzmann specific heat $C(q)$ provides a quantization criterion distinct from variance-based methods.
Results show \textbf{progressive cooling exists} in standard MHA models ($\rho=-0.66$, $p<0.001$) but is disrupted by GQA.
$C(q)$ and $\|k\|^2$ disagree on 39.8\% of bit allocations, confirming they measure different quantities, though $C(q)$ does not yet outperform variance-based quantization in PPL.
\end{abstract}

\section{Experiment 1-1: $T_\mathrm{eff}$ Gate Test}

\subsection{Setup}
The effective temperature is defined as $T_\mathrm{eff}^{(l)} = S^{(l)} / \log n$, where $S^{(l)} = -\sum_j \alpha_j \log \alpha_j$ is the attention entropy.
We compute $T_\mathrm{eff}$ per layer and head by extracting $Q, K$ via forward hooks, then computing attention weights $\alpha = \mathrm{softmax}(QK^\top / \sqrt{d})$ manually.

\begin{itemize}[nosep]
  \item \textbf{Models}: Qwen2.5-3B (GQA, 2 KV heads), GPT-2 Medium (standard MHA)
  \item \textbf{Calibration}: 64 samples, $n_\mathrm{max} = 512$
  \item \textbf{Gate metric}: Spearman $\rho(\text{layer index}, T_\mathrm{eff})$
\end{itemize}

\subsection{Results}

\begin{table}[H]
\centering
\caption{Experiment 1-1 gate results across two architectures.}
\label{tab:gate}
\begin{tabular}{lcccl}
\toprule
\textbf{Model} & \textbf{GQA} & $\boldsymbol{\rho}$ & $\boldsymbol{p}$ & \textbf{Gate} \\
\midrule
Qwen2.5-3B & 2 KV heads & $+0.108$ & $0.53$ & \textcolor{failure}{FAILURE} \\
GPT-2 Medium & None (MHA) & $-0.657$ & $4.8\times10^{-4}$ & \textcolor{partial}{WEAK SUCCESS} \\
\bottomrule
\end{tabular}
\end{table}

\paragraph{GPT-2 Medium (MHA).}
Clear progressive cooling: $T_\mathrm{eff}$ decreases from 0.86 (Layer~0) to 0.22 (Layer~19), with a final-layer ``reheating'' to 0.52 (Layer~23).
This pattern is consistent with the Boltzmann interpretation of multi-layer attention as progressive annealing.

\paragraph{Qwen2.5-3B (GQA).}
A V-shaped pattern: hot initial layers (0.84) $\to$ cold middle (0.30--0.60) $\to$ reheat (0.76).
No monotonic trend ($\rho = +0.11$).
The GQA mechanism introduces extreme inter-group temperature gaps---at Layer~5, KV-group~0 has $T_\mathrm{eff}=0.548$ while KV-group~1 has $T_\mathrm{eff}=0.042$, a gap of \textbf{0.506}.
Head variance in middle layers ($\sigma^2_h = 0.036$) is $7\times$ that of initial/final layers (0.005), reflecting strong functional differentiation between KV groups.

\paragraph{Implication.}
The Boltzmann thermal structure is a \textit{real phenomenon} in standard MHA transformers.
GQA disrupts it by forcing multiple query heads to share key-value projections, creating artificial temperature patterns.
Future work should analyze $T_\mathrm{eff}$ \textit{per KV group} rather than per head in GQA models.

\subsection{Gate Decision}
Per the pre-registered criteria: $\rho < -0.7$ (strong), $-0.7 < \rho < -0.3$ (weak), $\rho > -0.3$ (failure).
GPT-2 Medium achieves \textcolor{partial}{weak success} ($\rho = -0.66$, just above the strong threshold).
\textbf{Proceed to Experiments 1-2 and 1-3.}

\section{Experiment 1-2: Specific Heat Quantization}

\subsection{Setup}
The Boltzmann specific heat is $C(q_i) = \mathrm{Var}_\alpha[q_i^\top k_j / \sqrt{d}]$, the attention-weighted variance of energy.
We compare three KV cache quantization strategies at 2, 3, and 4-bit:

\begin{enumerate}[nosep]
  \item \textbf{Uniform}: all tokens quantized at the same bit width
  \item \textbf{Variance-based}: $\|k_j\|^2 > \text{median} \to +1$~bit, else $-1$~bit
  \item \textbf{Specific heat}: $C(q_i) > \text{median} \to +1$~bit, else $-1$~bit
\end{enumerate}

Model: GPT-2 Medium. Evaluation: WikiText-2 (32 samples, $n_\mathrm{max}=256$). Baseline PPL: 139.85.

\subsection{Disagreement Analysis}

\begin{table}[H]
\centering
\caption{Bit allocation disagreement between $C(q)$ and $\|k\|^2$ by layer group.}
\label{tab:disagree}
\begin{tabular}{lcc}
\toprule
\textbf{Layer Group} & \textbf{Avg Disagreement} & \textbf{Avg $\rho(C, \|k\|^2)$} \\
\midrule
Early (L0--L5) & 24.4\% & 0.56 \\
Middle (L6--L17) & 42.3\% & 0.19 \\
Late (L18--L23) & 42.6\% & 0.14 \\
\midrule
\textbf{Overall} & \textbf{39.8\%} & \textbf{0.275} \\
\bottomrule
\end{tabular}
\end{table}

$C(q)$ and $\|k\|^2$ disagree on \textbf{39.8\%} of bit allocations (well above the 20\% threshold).
Disagreement increases in later layers where attention is more focused, reaching 49.6\% at Layer~21.
The Spearman correlation between the two quantities is only 0.275, confirming they capture largely independent information.

\subsection{PPL Results}

\begin{table}[H]
\centering
\caption{Perplexity under different quantization strategies. Baseline (FP32): 139.85.}
\label{tab:ppl}
\begin{tabular}{lccc}
\toprule
\textbf{Bits} & \textbf{Uniform} & \textbf{Variance} & \textbf{Specific Heat} \\
\midrule
2-bit & 419.4 & $\infty$\textsuperscript{*} & $\infty$\textsuperscript{*} \\
3-bit & 130.4 & \textbf{127.6} & 264.5 \\
4-bit & \textbf{110.2} & 122.9 & 125.6 \\
\bottomrule
\multicolumn{4}{l}{\footnotesize \textsuperscript{*}Numerical instability in adaptive 2-bit scheme.}
\end{tabular}
\end{table}

\subsection{Decision: \textcolor{partial}{DIFFERENT BUT NOT BETTER}}

$C(q)$ provides a \textbf{genuinely novel measurement} that captures query-dependent sensitivity missed by key-norm methods.
However, the naive median-based $\pm1$~bit allocation does not translate this into PPL improvement.
Possible reasons:
\begin{itemize}[nosep]
  \item The median threshold is too coarse---$C(q)$'s value lies in identifying \textit{extreme} sensitivity, not binary classification.
  \item The hook-based K-replacement may introduce artifacts (2-bit explosion in both adaptive methods).
  \item $C(q)$ may need to be combined with $\|k\|^2$ rather than replacing it.
\end{itemize}

\section{Summary and Next Steps}

\begin{table}[H]
\centering
\caption{Level 1 decision matrix (current status).}
\begin{tabularx}{\textwidth}{lll>{\raggedright\arraybackslash}X}
\toprule
\textbf{Exp} & \textbf{Result} & \textbf{Gate} & \textbf{Implication} \\
\midrule
1-1 Gate & $\rho=-0.66$ & \textcolor{partial}{Weak Success} & Boltzmann thermal structure confirmed in MHA \\
1-2 Quant & 39.8\% disagree & \textcolor{partial}{Different, not better} & $C(q)$ is novel but needs refined allocation \\
1-3 FFN & --- & Pending & --- \\
\bottomrule
\end{tabularx}
\end{table}

\paragraph{Key contributions so far.}
\begin{enumerate}[nosep]
  \item Progressive cooling ($T_\mathrm{eff}$ monotonic decrease) is empirically confirmed in MHA.
  \item GQA disrupts this pattern---a novel finding with implications for GQA-aware analysis.
  \item $C(q)$ and $\|k\|^2$ are largely independent ($\rho=0.275$, 39.8\% disagreement), validating $C(q)$ as a \textit{new} quantity, not a repackaging.
  \item Translating $C(q)$ into practical quantization gain requires more sophisticated allocation.
\end{enumerate}

\paragraph{Immediate next steps.}
\begin{itemize}[nosep]
  \item \textbf{Exp 1-3}: Test whether FFN enhances spectral frequency selectivity.
  \item \textbf{Refined 1-2}: Try top-$k$\% / bottom-$k$\% allocation instead of median split; combine $C(q) + \|k\|^2$ as joint criterion.
  \item \textbf{Scale validation}: Test on larger non-GQA model (GPT-2 XL, 1.5B).
\end{itemize}

\end{document}
